% GENERATED FILE. Edit canonical lessons, then run npm run generate:books.
% canonical-source-hash: fnv1a64:f6c362da9af0dce4
% canonical-lessons: TE-C548-kacu, TE-C548-madcu, TE-C548-tarugu, TE-C548-jallincu, TE-C548-talimpu

\chapter{In the Kitchen}
\label{ch:te-in-the-kitchen}
\hlchaptermodality{548}

\begin{chapteropening}
\textbf{By the end of this chapter:} \emph{I can talk about heating, burning, chopping, sifting and tempering.}

Complete the last lesson of chapter 548: I can talk about heating, burning, chopping, sifting and tempering.
\end{chapteropening}

\section[\texorpdfstring{\te{కాచు} (kācu)}{కాచు (kācu)}]{\te{కాచు} (kācu) — to boil milk; to heat}

\label{lesson:TE-C548-kacu}



\begin{quote}
Before the new one: say the Telugu for to shake out, then the Telugu for to hang up.
\end{quote}

\subsection*{You'll want to know: \te{కాచు}}
\textbf{\te{కాచు}} — \emph{kācu} — \textquotedblleft{}to boil milk; to heat\textquotedblright{}.

\begin{quote}
\te{అమ్మ} \te{పొద్దున్నే} \te{పాలు} \te{కాచింది}.
\end{quote}

\begin{quote}
\emph{amma poddunnē pālu kācindi.}
\end{quote}

\begin{quote}
Mum boiled the milk early in the morning.
\end{quote}

Past (I …): \te{కాచాను}, \emph{kācānu}.

For milk and water; \te{అన్నం} \te{వండు} for rice.

\begin{cousinweb}
\textbf{\te{కాచు}} (\emph{kācu}) is native Dravidian: DEDR 1458 lists Telugu \emph{kācu}, \textquotedblleft{}to heat, to warm, to boil\textquotedblright{}, with Tamil \emph{kāy}, \textquotedblleft{}to grow hot\textquotedblright{}.
\end{cousinweb}

\begin{grammarlens}[title={What you've built}]
One more word for describing what you lived and felt.
\end{grammarlens}

\subsection*{Your turn}
\begin{itemize}
\raggedright
  \item \emph{Say it:} \emph{kācu}
  \item \emph{Say it:} \emph{kācu}, once more
  \item \emph{Say it:} \emph{amma poddunnē pālu kācindi.} — Mum boiled the milk early in the morning.
  \item \emph{Recall:} say the Telugu for to shake out, then the Telugu for to hang up, then say \emph{kācu} again
\end{itemize}

\begin{tcolorbox}[breakable,colback=teal!4,colframe=teal!35!black,title={Before you move on}]
What does \te{కాచు} mean? (\textquotedblleft{}To boil milk; to heat\textquotedblright{}.) Say it once more.
\end{tcolorbox}
\section[\texorpdfstring{\te{మాడ్చు} (māḍcu)}{మాడ్చు (māḍcu)}]{\te{మాడ్చు} (māḍcu) — to scorch; to burn food}

\label{lesson:TE-C548-madcu}



\begin{quote}
Before the new one: say the Telugu for to hang up, then the Telugu for to boil milk.
\end{quote}

\subsection*{You'll want to know: \te{మాడ్చు}}
\textbf{\te{మాడ్చు}} — \emph{māḍcu} — \textquotedblleft{}to scorch; to burn food\textquotedblright{}.

\begin{quote}
\te{ఫోనులో} \te{మాట్లాడుతూ} \te{నేను} \te{కూర} \te{మాడ్చాను}.
\end{quote}

\begin{quote}
\emph{phōnulō māṭlāḍutū nēnu kūra māḍcānu.}
\end{quote}

\begin{quote}
Talking on the phone, I burnt the curry.
\end{quote}

Past (I …): \te{మాడ్చాను}, \emph{māḍcānu}.

\te{మాడు}, to get burnt, is its intransitive partner.

\begin{cousinweb}
\textbf{\te{మాడ్చు}} (\emph{māḍcu}) is native Dravidian: DEDR 4801 lists Telugu \emph{māḍu}, \textquotedblleft{}to be burnt, to be scorched black\textquotedblright{}, and \emph{māḍ(u)cu}, \textquotedblleft{}to burn, to scorch\textquotedblright{}.
\end{cousinweb}

\begin{grammarlens}[title={What you've built}]
One more word for describing what you lived and felt.
\end{grammarlens}

\subsection*{Your turn}
\begin{itemize}
\raggedright
  \item \emph{Say it:} \emph{māḍcu}
  \item \emph{Say it:} \emph{māḍcu}, once more
  \item \emph{Say it:} \emph{phōnulō māṭlāḍutū nēnu kūra māḍcānu.} — Talking on the phone, I burnt the curry.
  \item \emph{Recall:} say the Telugu for to hang up, then the Telugu for to boil milk, then say \emph{māḍcu} again
\end{itemize}

\begin{tcolorbox}[breakable,colback=teal!4,colframe=teal!35!black,title={Before you move on}]
What does \te{మాడ్చు} mean? (\textquotedblleft{}To scorch; to burn food\textquotedblright{}.) Say it once more.
\end{tcolorbox}
\section[\texorpdfstring{\te{తరుగు} (tarugu)}{తరుగు (tarugu)}]{\te{తరుగు} (tarugu) — to chop; to shred}

\label{lesson:TE-C548-tarugu}



\begin{quote}
Before the new one: say the Telugu for to boil milk, then the Telugu for to scorch.
\end{quote}

\subsection*{You'll want to know: \te{తరుగు}}
\textbf{\te{తరుగు}} — \emph{tarugu} — \textquotedblleft{}to chop; to shred\textquotedblright{}.

\begin{quote}
\te{వంట} \te{కోసం} \te{నేను} \te{ఉల్లిపాయలు} \te{తరిగాను}.
\end{quote}

\begin{quote}
\emph{vaṇṭa kōsaṁ nēnu ullipāyalu tarigānu.}
\end{quote}

\begin{quote}
I chopped the onions for the cooking.
\end{quote}

Past (I …): \te{తరిగాను}, \emph{tarigānu}.

Fine cutting of vegetables; \te{కోయు} is the general word for cutting.

\begin{cousinweb}
\textbf{\te{తరుగు}} (\emph{tarugu}) is native Dravidian: DEDR 3140 lists Telugu \emph{taṟugu}, \textquotedblleft{}to slice, to chop\textquotedblright{}, with Tamil \emph{taṟi}, \textquotedblleft{}to cut off\textquotedblright{}.
\end{cousinweb}

\begin{grammarlens}[title={What you've built}]
One more word for describing what you lived and felt.
\end{grammarlens}

\subsection*{Your turn}
\begin{itemize}
\raggedright
  \item \emph{Say it:} \emph{tarugu}
  \item \emph{Say it:} \emph{tarugu}, once more
  \item \emph{Say it:} \emph{vaṇṭa kōsaṁ nēnu ullipāyalu tarigānu.} — I chopped the onions for the cooking.
  \item \emph{Recall:} say the Telugu for to boil milk, then the Telugu for to scorch, then say \emph{tarugu} again
\end{itemize}

\begin{tcolorbox}[breakable,colback=teal!4,colframe=teal!35!black,title={Before you move on}]
What does \te{తరుగు} mean? (\textquotedblleft{}To chop; to shred\textquotedblright{}.) Say it once more.
\end{tcolorbox}
\section[\texorpdfstring{\te{జల్లించు} (jalliñcu)}{జల్లించు (jalliñcu)}]{\te{జల్లించు} (jalliñcu) — to sift}

\label{lesson:TE-C548-jallincu}



\begin{quote}
Before the new one: say the Telugu for to scorch, then the Telugu for to chop.
\end{quote}

\subsection*{You'll want to know: \te{జల్లించు}}
\textbf{\te{జల్లించు}} — \emph{jalliñcu} — \textquotedblleft{}to sift\textquotedblright{}.

\begin{quote}
\te{అమ్మ} \te{పిండిని} \te{జల్లించింది}.
\end{quote}

\begin{quote}
\emph{amma piṇḍini jalliñcindi.}
\end{quote}

\begin{quote}
Mum sifted the flour.
\end{quote}

Past (I …): \te{జల్లించాను}, \emph{jalliñcānu}.

Flour, rice or sand is sifted through a \te{జల్లెడ}, a sieve, to take out the lumps and grit.

\begin{cousinweb}
\textbf{\te{జల్లించు}} (\emph{jalliñcu}) is native Dravidian: DEDR 2370 lists Telugu \emph{jallincu}, \textquotedblleft{}to sift\textquotedblright{}, and \emph{jalleḍa}, \textquotedblleft{}a sieve\textquotedblright{}, with Tamil \emph{cali}.
\end{cousinweb}

\begin{grammarlens}[title={What you've built}]
One more word for describing what you lived and felt.
\end{grammarlens}

\subsection*{Your turn}
\begin{itemize}
\raggedright
  \item \emph{Say it:} \emph{jalliñcu}
  \item \emph{Say it:} \emph{jalliñcu}, once more
  \item \emph{Say it:} \emph{amma piṇḍini jalliñcindi.} — Mum sifted the flour.
  \item \emph{Recall:} say the Telugu for to scorch, then the Telugu for to chop, then say \emph{jalliñcu} again
\end{itemize}

\begin{tcolorbox}[breakable,colback=teal!4,colframe=teal!35!black,title={Before you move on}]
What does \te{జల్లించు} mean? (\textquotedblleft{}To sift\textquotedblright{}.) Say it once more.
\end{tcolorbox}
\section[\texorpdfstring{\te{తాలింపు} (tālimpu)}{తాలింపు (tālimpu)}]{\te{తాలింపు} (tālimpu) — tempering; seasoning}

\label{lesson:TE-C548-talimpu}



\begin{quote}
Before the new one: say the Telugu for to chop, then the Telugu for to sift.
\end{quote}

\subsection*{You'll want to know: \te{తాలింపు}}
\textbf{\te{తాలింపు}} — \emph{tālimpu} — \textquotedblleft{}tempering; seasoning\textquotedblright{}.

\begin{quote}
\te{ఆవాలు}, \te{కరివేపాకుతో} \te{పప్పుకు} \te{తాలింపు} \te{పెట్టారు}.
\end{quote}

\begin{quote}
\emph{āvālu, karivēpākutō pappuku tālimpu peṭṭāru.}
\end{quote}

\begin{quote}
They tempered the dal with mustard seeds and curry leaves.
\end{quote}

Spices fried in oil and poured over a dish: \te{తాలింపు} \te{వేయు}.

\begin{cousinweb}
\textbf{\te{తాలింపు}} (\emph{tālimpu}) is native Dravidian: DEDR 3186 lists Telugu \emph{tālimpu}, \textquotedblleft{}seasoning\textquotedblright{}, and \emph{tālincu}, \textquotedblleft{}to season\textquotedblright{}, with Tamil \emph{tāḷi}, \textquotedblleft{}to season with spices fried in ghee\textquotedblright{}.
\end{cousinweb}

\begin{grammarlens}[title={What you've built}]
One more word for describing what you lived and felt.
\end{grammarlens}

\subsection*{Your turn}
\begin{itemize}
\raggedright
  \item \emph{Say it:} \emph{tālimpu}
  \item \emph{Say it:} \emph{tālimpu}, once more
  \item \emph{Say it:} \emph{āvālu, karivēpākutō pappuku tālimpu peṭṭāru.} — They tempered the dal with mustard seeds and curry leaves.
  \item \emph{Recall:} say the Telugu for to chop, then the Telugu for to sift, then say \emph{tālimpu} again
\end{itemize}

\begin{tcolorbox}[breakable,colback=teal!4,colframe=teal!35!black,title={Before you move on}]
What does \te{తాలింపు} mean? (\textquotedblleft{}Tempering; seasoning\textquotedblright{}.) Say it once more.
\end{tcolorbox}
